导体中存在能够在材料内部宏观移动的自由电荷；绝缘体中的电荷则主要束缚在原子、分子或晶格
附近。外电场加上以后，导体内的自由电荷会重新分布，直到不再发生定向移动。此时虽然电荷已经
静止，表面电荷产生的附加电场却仍然存在，并与外电场共同构成最终的静电场。

\subsection*{导体的静电平衡}

设导体已经达到静电平衡。若导体内部仍有非零电场，自由电荷便会继续受力移动，与“平衡”矛盾，
所以
\begin{equation}
  \boxed{\vect E_{\rm in}=0}.
\end{equation}
由 $\vect E=-\nabla\varphi$ 可知，整块连通导体是等势体，其表面也是等势面。导体内部任取高斯面，
因高斯面上电场处处为零，包围的净电荷必为零。因此没有置于空腔中的独立电荷时，导体的过剩
净电荷只能分布在表面。

表面若存在切向电场，自由电荷会沿表面移动；故平衡时表面外的电场只能沿法线方向。取跨越导体
表面的薄柱形高斯面，可得
\begin{equation}
  \boxed{E_n^{\rm(out)}-E_n^{\rm(in)}=\frac{\sigma}{\varepsilon_0}}.
\end{equation}
导体内侧 $E_n^{\rm(in)}=0$，于是紧邻表面外侧有
\begin{equation}
  \boxed{E_n^{\rm(out)}=\frac{\sigma}{\varepsilon_0}}.
\end{equation}
这里的正法线由导体指向外部；$\sigma>0$ 时电场向外，$\sigma<0$ 时电场向内。该关系是局部
关系，不要求导体为球形，也不能把弯曲表面的电场误写成无限大平面两侧的
$\sigma/(2\varepsilon_0)$。

\twoConductorFieldDiagram

\conductorEquilibriumDiagram

把外电荷靠近中性导体时，导体两侧出现异号感应电荷，但导体的总电荷仍为零。若要用接地得到
带电导体，操作次序不能颠倒：保持外电荷不动时接地，使一部分自由电荷经大地流入或流出；随后
先断开地线，再移走外电荷。若先移走外电荷而仍保持接地，导体会重新与大地交换电荷，通常回到
中性状态。

\inductionGroundingDiagram

闭合导体空腔内若没有放置电荷，腔内边界电势为常量，唯一的静电解就是常量电势，因而腔内
$\vect E=0$。这就是静电屏蔽。它并不是说导体“阻挡”了已有的电场线，而是表面感应电荷重新
分布后，使空腔中满足边界条件的唯一解恰好为零场。

\historyfigure[0.42\linewidth]{faraday-cage-redraw.png}{Faraday 笼：放电止于导电笼外表面}

图中应观察的是长电弧终止在导电笼的外表面，而笼内的人并不处在放电通道中。外加电场使
笼上的自由电荷重新分布，在静电或变化足够缓慢的情况下，导体内部保持近似等势，笼内场强
很小。若电场快速变化，实际屏蔽效果还与网孔尺寸、材料和频率有关；因此实验图体现的是
静电屏蔽的边界条件，而不是“任何金属网在任何情况下都能完全隔绝电场”。

\subsection*{静电场唯一性定理}

电势满足 Poisson 方程
\begin{equation}
  \boxed{\nabla^2\varphi=-\frac{\rho}{\varepsilon_0}},
\end{equation}
无自由电荷区域则满足 Laplace 方程 $\nabla^2\varphi=0$。给定区域内的电荷密度，并在边界上
给定电势；或给定适当的法向电场且固定一个电势参考后，区域内的静电势由这些条件唯一确定。

以给定边界电势为例，假设 $\varphi_1$、$\varphi_2$ 都是解，令
$u=\varphi_1-\varphi_2$。则 $\nabla^2u=0$ 且边界上 $u=0$。由 Green 恒等式
\begin{equation}
  \int_V|\nabla u|^2\dd V
  =\oint_{\partial V}u\frac{\partial u}{\partial n}\dd S
   -\int_Vu\nabla^2u\dd V=0.
\end{equation}
因此 $\nabla u=0$，而边界上的零值又使 $u=0$，两组解只能相同。这个结论使我们不必直接追踪
导体表面连续分布的感应电荷：只要用若干辅助点电荷构造出的电势满足相同的方程和边界条件，
它就是所求区域中的真实电势。

\discretePotentialDiagram

\subsection*{电像法}

半径为 $R$ 的接地导体球，球心为 $O$。在球外距球心 $L>R$ 处放置点电荷 $Q$。在球内同一直线
上放置辅助电荷
\begin{equation}
  \boxed{q'=-\frac{R}{L}Q,\qquad l=\frac{R^2}{L}},
\end{equation}
其中 $l$ 是像电荷到球心的距离。球面上任一点 $P$ 到 $Q$ 和 $q'$ 的距离之比恒为 $L/R$，故
\begin{equation}
  \left.\varphi(P)\right|_{r=R}
  =k\left(\frac{Q}{PQ}+\frac{q'}{Pq'}\right)=0,
  \qquad k=\frac1{4\pi\varepsilon_0}.
\end{equation}
球外真实区域内没有像电荷，且球面电势与接地条件一致；由唯一性定理，这个“真实电荷加像电荷”
的场就是球外的真实电场。

\sphereImageDiagram

像电荷只是求导体外场的数学替代物，不能解释成导体内部真的聚集了一个点电荷。真实点电荷受到
导体的力，却可以直接按它与像电荷之间的库仑力计算：
\begin{equation}
  \boxed{F=-k\frac{Q^2RL}{(L^2-R^2)^2}},
\end{equation}
负号表示力指向球心。若 $\theta$ 是球面点的半径与 $OQ$ 的夹角，球面的感应电荷密度为
\begin{equation}
  \boxed{\sigma(\theta)=
  -\frac{Q(L^2-R^2)}{4\pi R
  \left(L^2+R^2-2LR\cos\theta\right)^{3/2}}}.
\end{equation}
积分可得接地球从大地取得的总感应电荷 $Q_{\rm ind}=-QR/L$，恰等于像电荷的电荷量。

力的计算不需要附加因子 $1/2$；系统能量却需要。导体的感应电荷随真实电荷逐渐移入而建立，故
接地球与外点电荷系统的能量为
\begin{equation}
  \boxed{U=\frac12Q\varphi_{\rm ind}(Q)
  =-\frac{kQ^2R}{2(L^2-R^2)}}.
\end{equation}

\subsection*{等离子体振荡}

等离子体中正、负电荷在宏观上近似中和。设电子数密度为 $n$，把电子云相对于静止的正离子背景
整体平移一个很小的距离 $\xi$。两端留下等效面电荷密度 $\pm ne\xi$，因而内部回复电场大小为
\begin{equation}
  E=\frac{ne\xi}{\varepsilon_0}.
\end{equation}
单个电子满足
\begin{equation}
  m\ddot\xi=-eE=-\frac{ne^2}{\varepsilon_0}\xi,
\end{equation}
所以电子云作简谐振荡，其等离子体角频率为
\begin{equation}
  \boxed{\omega_p=\sqrt{\frac{ne^2}{m\varepsilon_0}}}.
\end{equation}
更一般地，若载流粒子的质量为 $m$、电荷量绝对值为 $|q|$，则
$\omega_p=\sqrt{nq^2/(m\varepsilon_0)}$。该推导忽略了热运动、碰撞、边缘效应和离子运动，适合
描述均匀等离子体的长波小振幅振荡。

\plasmaOscillationDiagram

\ionosphereDiagram

\subsection*{导体边界上的局部规律}

静电平衡时，导体表面电场沿法线，但不同位置的场强一般不同。由
$E_n=\sigma/\varepsilon_0$ 可知，表面曲率较大、尖锐程度较高的位置往往需要更高的面电荷密度
来维持整个导体等势，因而附近场强更大。这是尖端放电和避雷针附近容易发生气体电离的基础。
“电荷集中在尖端”不是额外定律，而是等势边界条件与几何共同决定的结果。

在导体表面取一个电荷元 $\sigma\dd S$，它不能对自身施力。表面外侧总场为
$E_{\rm out}=\sigma/\varepsilon_0$，内侧为零；去除该电荷元自身在两侧方向相反的场后，其他电荷
对它的平均场为
\begin{equation}
 E_{\rm others}=\frac{E_{\rm out}+E_{\rm in}}2
 =\frac\sigma{2\varepsilon_0}.
\end{equation}
因此导体表面受到向外的静电压强
\begin{equation}
 \boxed{p=\sigma E_{\rm others}=\frac{\sigma^2}{2\varepsilon_0}
 =\frac12\varepsilon_0E_{\rm out}^2}.
\end{equation}
它与真空电场能量密度相等，可用于求带电球壳的膨胀趋势和电容器极板间的吸力。

\subsection*{空腔、有内电荷与接地}

若闭合导体空腔内放入电荷 $q$，在包围空腔而仍处于导体材料内部的高斯面上 $\vect E=0$，故
空腔内壁的总感应电荷必须为 $-q$。若导体原来总电荷为 $Q$，外表面的总电荷便为 $Q+q$。
这些结论只确定各表面的总量；具体面密度仍由空腔形状、内电荷位置和外部边界条件决定。

内电荷会使空腔中出现电场，却不会破坏导体材料内部的零场。外部场在只给定外表面总电荷和外部
边界以后，也不需要知道空腔中电荷的具体位置。这是静电屏蔽的另一层含义：闭合导体把内外两个
边值问题分开。若导体接地，外表面还可与大地交换任意电荷，外部场甚至可能完全不显露腔内电荷。

接地条件不是“导体电荷为零”，而是导体电势固定为大地电势，通常取零。接地导体可以带有非零
净电荷；孤立中性导体则必须保持总电荷为零，但其电势一般不为零。电像法中若把接地球改为孤立
中性球，原来的单个像电荷只能满足球面等势，不能同时满足总感应电荷为零；还需在球心再加入一个
辅助像电荷，使所有像电荷代数和为零。

\subsection*{唯一性定理的使用步骤}

用唯一性定理解决导体问题时，应依次检查：所研究区域内的真实体电荷是否相同；每个边界上的
电势或法向场条件是否相同；无穷远条件是否相同。只要三者满足，构造解便是该区域的唯一解。
像电荷可以放在研究区域之外，例如导体材料内部，因为那里并不要求构造场等于真实场；若把像
电荷误放入要求求解的真空区域，就改变了 Poisson 方程的电荷源，不能再用唯一性定理证明等价。

接地平面外点电荷是最简单的例子。把平面取为 $z=0$，真实电荷 $Q$ 位于 $z=a$，在 $z=-a$
放置像电荷 $-Q$。上半空间电势
\begin{equation}
 \varphi=\frac1{4\pi\varepsilon_0}
 \left(\frac Q{r_+}-\frac Q{r_-}\right)
\end{equation}
在平面上处处为零，且上半空间只有真实电荷，因此就是所求解。真实电荷受力等于它与像电荷间的
库仑力：
\begin{equation}
 \boxed{F_z=-\frac{Q^2}{16\pi\varepsilon_0a^2}}.
\end{equation}
导体表面感应电荷密度的积分为 $-Q$。这并不意味着一块有限孤立金属板必然感应出总量 $-Q$；
该结论依赖无限接地平面的边界条件。

\subsection*{像电荷能量中的二分之一}

求真实电荷受到的瞬时力时，导体当时已经完成感应，故可把像电荷当作固定场源直接使用，不乘
$1/2$。求系统能量时，感应电荷却是随着真实电荷从无穷远移入而逐渐建立的。以接地平面为例，
像电荷与真实电荷相距 $2a$，若直接写两点电荷相互作用能会得到
$-kQ^2/(2a)$；真实系统能量应为其一半：
\begin{equation}
 U=-\frac{kQ^2}{4a}=-\frac{Q^2}{16\pi\varepsilon_0a}.
\end{equation}
对 $a$ 求负导数恰好恢复上面的吸力。力和能量两种方法给出同一结果，是检查二分之一因子是否
处理正确的可靠办法。

\subsection*{等离子体振荡的物理尺度}

$\omega_p$ 随数密度的平方根增大，随载流子质量的平方根减小。电子比离子轻得多，所以高频响应
通常由电子主导，而离子在一个电子振荡周期内近似静止。当外加电场频率低于或接近等离子体频率
时，电子有足够时间重排并显著屏蔽外场；频率远高于 $\omega_p$ 时，电子来不及跟随，介质的响应
性质会改变。

小位移推导假设电子云整体平移，回复力与位移成正比。若振幅太大、密度明显不均匀、碰撞频繁或
存在磁场，振荡会出现阻尼、色散和不同模式，单一的简谐频率只给出最基本的数量级。尽管如此，
$\omega_p=\sqrt{ne^2/(m\varepsilon_0)}$ 仍是判断等离子体时间尺度的首要参数。

\subsection*{静电平衡、空腔与屏蔽}

导体内有可自由移动的电荷。若内部存在非零电场，自由电荷会继续运动，故静电平衡时导体内部
$\vec E=0$，整个导体及其表面为等势体。表面切向电场必须为零，否则表面电荷也会运动；电场只能
沿法线。用跨过表面的薄柱高斯面可得外侧法向场
\[
 E_n=\frac\sigma{\varepsilon_0}.
\]
多余净电荷分布在外表面，但分布通常不均匀：几何越尖锐处，为维持等势所需的面电荷密度越大，
邻近场强也越强。

导体空腔内若没有自由电荷，腔壁又是同一等势面，唯一性定理保证腔内电势为常量、场为零，外界
电场不能穿入，这就是静电屏蔽。若腔内放入电荷 $q$，腔壁总感应电荷为 $-q$；导体原有净电荷与
$q$ 的影响共同决定外表面总电荷。腔内场受腔内电荷和腔形决定，而在封闭导体中移动腔内电荷，
不会改变由外部边界条件唯一确定的外部场。

唯一性定理的内容是：给定区域内电荷分布，并在边界上给定电势或法向电场后，满足 Poisson 或
Laplace 方程的解至多一个。证明可设有两个解，考察其差；差函数满足 Laplace 方程且边界为零，
由能量积分或极值原理只能处处为零。此定理让我们不必直接解偏微分方程：只要构造一个满足场方程
和全部边界条件的候选解，它就是物理解。

\subsection*{电像法与受力、能量}

电像法用导体区域以外的虚拟电荷替代未知表面电荷，使导体边界成为规定等势面。接地无限平面前
的点电荷 $q$，可在镜像位置放置 $-q$；真实半空间中的电势与场便完全正确。像电荷不是真实电荷，
不能在真实空间里另算其电荷守恒，也不能求“像电荷受到的力”。真实电荷所受力可由像电荷产生的
场计算。

能量计算要特别注意装配过程。点电荷靠近平面时，像电荷的位置会随真实电荷移动，系统能量不是
简单的真实电荷与像电荷相互作用能，而是其一半：
\[
 U=\frac12q\varphi_{\rm image}(\vec r_q).
\]
对固定位置直接求力则不需要这个 $1/2$。球面边界的像电荷位置和电量由反演关系决定；接地球与
孤立中性球的边界条件不同，后者还要增加中心像电荷以保证导体总电荷为零。

等离子体中电子相对几乎不动的离子背景发生微小整体位移时，会出现把电子拉回平衡位置的电场。
电子密度为 $n$ 时，线性化运动方程给出等离子体频率
\[
 \omega_p=\sqrt{\frac{ne^2}{m_e\varepsilon_0}}.
\]
该频率与振幅无关是小振动近似的结果；它说明导体屏蔽并非“电荷瞬时到位”，自由载流子有自身的
集体响应时间尺度。


\subsection*{从微观移动到宏观边界条件}

导体与绝缘体的区别，不在于材料中“有没有电荷”，而在于是否存在能在所研究时间尺度内作宏观
迁移的载流子。金属中的自由电子只要受到很小的切向电场，就会产生定向运动；因此静电平衡不是
外电场突然消失，而是载流子经过一个暂态重排后，表面电荷产生的附加电场恰好把导体内部的总场
抵消。绝缘体中的束缚电荷也会发生微小位移或取向极化，但通常不能像金属自由电子那样在整块材料
中重新分配。

由此可把导体静电平衡的结论串成一条逻辑链：导体材料内部 $\vect E=0$；于是导体内部电势没有
空间变化；连通导体连同表面成为一个等势体；表面切向场为零，外侧场只能沿法线；再用跨过表面
的薄柱高斯面，才得到 $E_n^{\rm(out)}=\sigma/\varepsilon_0$。这些结论必须按此顺序理解，不能
把“等势”误解成“表面没有电荷”，也不能把局部法向场直接套成球对称场。

图 31.1(a) 中，两块导体各自等势，但二者电势可以不同。实线从高电势导体出发并沿法线离开表面，
进入低电势导体时仍与表面正交；虚线表示等势线，处处与电场线正交。两导体之间的电场是一个
整体边值问题，不能把每块导体都当成互不影响的孤立带电球。表面某一点电荷密度的正负，应由该点
外侧场的法向方向判断，而不能只由整块导体的净电荷符号判断。

\subsection*{感应起电的全过程}

以带正电的外物体靠近中性金属球为例。最初，电子向近端移动，近端带负电、远端带等量正电，
金属球总电荷仍为零。保持外电荷位置不变并接地后，大地与球构成一个巨大的连通导体；在外正
电荷吸引下，电子从大地流入球体，原先远端的正电荷被中和。此时球的电势被固定为大地电势，
并不要求球的净电荷等于零。

随后必须先断开地线，使已经流入球内的电子失去返回大地的通道；最后再移去外电荷。外场撤除后，
球面上的电子在相互排斥下重新分布。若球是孤立球面且周围没有其他物体，最终分布具有球对称性；
若导体形状不规则或附近仍有其他导体，最终电荷仍会不均匀。“均匀分布”是球对称边界条件的结果，
不是感应起电对任意导体都成立的结论。若把最后两步颠倒，在仍然接地时先撤去外电荷，多余电子会
流回大地，球通常重新变为中性。

\subsection*{空腔与静电屏蔽的边界}

闭合导体空腔内没有自由电荷时，腔内满足 Laplace 方程，而整个腔壁属于同一连通导体、具有同一
电势。常量电势显然满足方程与边界条件；唯一性定理保证它就是唯一解，所以腔内电场为零。这个
结论并不需要腔体是球形，也不需要外界电场具有任何对称性。

若腔内放入电荷 $q$，腔内不再是零场。在导体材料中作包围空腔的高斯面，由于该面上场为零，
腔内壁的总感应电荷必为 $-q$。若孤立导体原有净电荷 $Q$，则外表面的总电荷为 $Q+q$；若导体
接地，外表面总电荷由零电势边界和外界环境决定。这里“屏蔽”表示内外边值问题受导体等势边界
约束，并不表示金属壁把电场线当作实体挡住。

真实 Faraday 笼只在静电或足够低频的近似下可直接用上述结论。时变场会在导体中激发电流，
屏蔽效果取决于电导率、壁厚、趋肤深度、网孔尺寸和电磁波频率。图中的放电通道终止在笼外表面，
说明电荷主要沿导体外部重新分布；它不能推出任意金属网对任意频率都能实现完全屏蔽。

\subsection*{从 Laplace 方程到离散电势}

一维无电荷区域满足 $\dd^2\varphi/\dd x^2=0$。把相邻网格点间距记为 $a$，在 $x_n$ 两侧作
Taylor 展开：
\begin{align*}
 \varphi_{n+1}&=\varphi_n+a\varphi'_n+\frac{a^2}{2}\varphi''_n+O(a^3),\\
 \varphi_{n-1}&=\varphi_n-a\varphi'_n+\frac{a^2}{2}\varphi''_n+O(a^3).
\end{align*}
两式相加，奇次项相消，于是
\begin{equation}
 \varphi''(x_n)=\frac{\varphi_{n+1}-2\varphi_n+\varphi_{n-1}}{a^2}+O(a^2).
\end{equation}
在二维等间距方格上，$\nabla^2\varphi=0$ 相应变成
\begin{equation}
 \boxed{\varphi_{i,j}=\frac14\bigl(\varphi_{i+1,j}+\varphi_{i-1,j}
 +\varphi_{i,j+1}+\varphi_{i,j-1}\bigr)}.
\end{equation}
也就是说，无电荷区域内部一点的电势等于四个最近邻点电势的平均值。这个“平均值性质”解释了
为什么调和函数不能在内部产生孤立的最大值或最小值，也给出逐点迭代求复杂导体电势的数值方法。
图 31.4(b) 中标出的中心二阶差分正是连续 Laplace 方程与离散电路、数值网格之间的桥梁。

唯一性定理在此还提供了验算标准：候选电势不仅要让导体表面“看起来像等势面”，还必须在整个
求解区域满足正确的 Poisson/Laplace 方程，并满足所有边界以及无穷远条件。只检验其中一部分，
不能证明构造解就是物理解。

\subsection*{球面像电荷的几何关系}

接地球问题中，真实电荷 $Q$ 位于距球心 $L$ 处，像电荷放在同一半径上、距球心
$l=R^2/L$ 处。任取球面点 $P$，由反演几何可得
\begin{equation}
 \frac{PQ}{Pq'}=\frac{L}{R}.
\end{equation}
选择 $q'=-(R/L)Q$ 后，$P$ 点两项电势恰好抵消，因而整个球面而非某一个特殊点都保持零电势。
这也是图 31.3 中像电荷必须位于球内且靠近真实电荷一侧的原因。若图中把像电荷画在错误位置，
就会破坏上述恒定距离比，边界条件不再成立。

像电荷法只保证所选真实区域内的电势正确。球内的像电荷不是金属中真实存在的点电荷；真实感应
电荷连续分布在球面，其近 $Q$ 一侧密度较大。计算 $Q$ 所受力时，可直接用像电荷的库仑力，因为
此时只需感应场在 $Q$ 位置的瞬时值。计算装配能量时必须写成
$U=Q\varphi_{\rm ind}/2$，因为感应电荷是在 $Q$ 从无穷远移近的过程中逐渐建立的。用
$F=-\dd U/\dd L$ 应当能恢复同一个吸引力，这是一项重要自检。

接地球和孤立球是不同边值问题。接地球可以与大地交换电荷，只要求球面电势为零；孤立中性球
既要等势，又要保证表面总感应电荷为零，通常需在球心补一个像电荷。若孤立球本来带给定净电荷，
球心像电荷的大小还要相应改变。因此看到“导体球”时，必须先判断它是否接地、是否孤立以及净
电荷是否给定，再选择像电荷组。

\subsection*{等离子体频率与电离层传播}

电子云相对离子背景移动距离 $\xi$ 后，两端出现的等效面电荷产生回复场，在线性近似下得到
$\omega_p=\sqrt{ne^2/(m\varepsilon_0)}$。这一频率描述自由载流子集体重排的固有时间尺度：
密度越高，回复场越强；质量越大，响应越慢。电子质量远小于离子质量，所以通常先讨论电子
等离子体振荡。

对简化的无碰撞等离子体，横向电磁波的色散关系可写为
\begin{equation}
 \omega^2=\omega_p^2+c^2k^2.
\end{equation}
当 $\omega<\omega_p$ 时，传播常数成为虚数，波在等离子体中迅速衰减，宏观上表现为反射；当
$\omega>\omega_p$ 时，波可以传播。电离层中的电子密度随高度、昼夜和太阳活动变化，因此反射
并不发生在一张固定的“镜面”上。图 31.6(a) 用弯曲路径表示低频无线电波可返回地面、支持远距离
通信，同时以穿透路径提醒读者：更高频率的波可能越过电离层。实际传播还会受到碰撞、地磁场、
分层结构和入射角的影响，等离子体频率只是判断能否传播的第一尺度。

\subsection*{导体系统的电势与电荷关系}

有多块彼此分离的导体时，每块导体内部仍是等势体，但各导体电势一般不同。在线性静电范围内，
导体电荷和电势可写成
\[
 Q_i=\sum_j C_{ij}\varphi_j,
\]
其中系数只由几何形状、相对位置和介质决定。对角元为正，非对角元通常为负；改变一块导体电势
会通过感应改变其他导体的表面电荷分布。这一形式提醒我们：两导体的净电荷大小不能单独决定谁的
电势更高，导体尺寸和相互位置同样重要。

判断表面是否可能出现异号电荷时，可结合最高、最低电势原理。在无自由体电荷区域，调和函数的
极值只能出现在边界。电势最高的导体表面，外法向电场不能指向导体内部，因而其表面不会出现
负电荷区；电势最低者相应不会出现正电荷区。但一个净电荷为正、却不是最高电势的导体，面对更
高电势导体的一侧仍可能出现局部负电荷。局部符号与净电荷符号并非同一问题。

\subsection*{边值题的检查清单}

构造静电解时，应逐项核对以下内容：真实求解区域内的电荷源是否正确；每一块接地导体是否为零
电势；每一块孤立导体是否等势且总电荷满足题设；不同边界相交处是否同时满足条件；无穷远处的
电势或场是否具有正确衰减。像电荷的数量少并不意味着解更“物理”，只有全部条件同时满足才可
借唯一性定理确认。

对两个相互垂直的接地平面，真实电荷关于每个平面各产生一个异号像，再对两个平面同时反射产生
一个同号像。四个电荷组成的势在两平面上分别成对抵消。若漏掉双重反射的同号像，其中一个平面
可能仍为零势，另一个却不满足边界。半球凸起或球壳内电荷问题同理：应先明确真实区域，再把像
电荷全部放到真实区域之外，并逐个验证平面反射和球面反演关系。

图中位置关系可用无量纲比值复核。接地球外点电荷必须满足
\[
 \frac{l}{R}=\frac{R}{L}<1,
\]
所以像电荷必在球内；球壳内点电荷 $x<R$ 的反演像位于 $R^2/x>R$，必在球外。只要图把两者画
反，即使公式文字正确，读者也会得到相反的空间判断。

\subsection*{数值解与物理解的对应}

复杂导体形状往往没有简单解析像电荷。把求解区域划分为网格，固定边界节点电势，再反复用
“内部节点等于邻点平均值”更新，就可逐步逼近 Laplace 方程的解。迭代收敛后，电场由相邻节点
电势差近似得到，导体表面电荷密度再由法向场
$\sigma=\varepsilon_0E_n$ 求出。网格越细，几何边界的表示越准确，但计算量也越大；尖角附近
场变化剧烈，通常需要局部加密。

数值图同样要接受物理检查：等势线不得相交；电场线应与导体表面正交；无体电荷区域不能出现
孤立电势极值；由包围导体的高斯面计算出的通量应与总电荷一致。数值收敛只说明离散方程被满足，
若边界条件输入错误，仍会稳定收敛到错误答案。

\subsection*{场能、静电压强与机械平衡}

导体表面外侧电场为 $E=\sigma/\varepsilon_0$，但某个表面电荷元受到的不是这个包含自场的总场。
去除自场后，其他电荷在表面处提供平均场 $E_{\rm other}=E/2$，因此表面受到向外的静电压强
\[
 p=\frac{\sigma^2}{2\varepsilon_0}
 =\frac12\varepsilon_0E^2.
\]
这个结果与真空电场能量密度相同。带电薄球壳有向外膨胀趋势，电容器两极板则相互吸引；方向
应结合两侧场强和边界法线判断，不能只看 $\sigma^2$ 的正值。

若导体由电源维持固定电势，机械位移时电源会输送或吸收电荷，单看电场能变化不足以给出机械力，
还要计入电源做功。孤立导体保持电荷不变，接地导体保持电势不变，这两类约束看似只差一句话，
却可能使能量法中的控制变量完全不同。先写清外界保持的量，再对几何参数求导，是避免符号错误
的关键。


\Needspace{0.34\textheight}
\begin{origexamplebox}{例题 31.1【异号导体】}
导体 $A$、$B$ 分别带电 $+Q$、$-q$，且 $Q>q$。由于二者会产生静电感应，问：导体 $A$
上是否可能出现负电荷？导体 $B$ 上是否可能出现正电荷？若 $Q=q$，结果有无变化？
\end{origexamplebox}

\begin{xesolution}
表面电荷的局部符号不能只看净电荷，而要结合导体电势。电场线从较高电势出发，终止于较低
电势或无穷远；同一导体是等势体，电场线不可能从它的一处出发又终止于它的另一处。

本题中 $V_A>V_B$。又因总电荷 $Q-q>0$，可知 $V_A>0$，所以 $A$ 是 $A$、$B$ 与无穷远三者
中的最高电势边界。最高电势导体表面外的电场只能指向外部，故
\[
  \boxed{A\text{ 的表面不可能出现负电荷区}}.
\]
$B$ 的电势可能为正、为零或为负，这取决于两导体的形状、大小和间距。当 $V_B>0$ 时，$B$ 的
背离 $A$ 一侧可以有少量正电荷，而面对 $A$ 的一侧有较多负电荷，总电荷仍为 $-q$。因此
\[
  \boxed{B\text{ 上可能出现正电荷，但不是必然出现}}.
\]

若 $Q=q$，总电荷为零。此时 $V_A>0>V_B$：$A$ 是最高电势边界，$B$ 是最低电势边界。于是
$A$ 表面处处非负，$B$ 表面处处非正；除去电荷密度恰为零的特殊点，可分别说成只带正电和只带
负电。因而 $Q=q$ 时，$B$ 上也不可能出现正电荷区。

\pairPotentialDiagram
\end{xesolution}

\Needspace{0.34\textheight}
\begin{origexamplebox}{例题 31.2【同号导体】}
导体 $A$、$B$ 分别带电 $+Q$、$+q$，且 $Q>q$。由于二者会产生静电感应，问：两导体上是否
可能出现负电荷？若 $Q=q$，结果有无变化？若有更多带正电的导体，例如 $100$ 个，能得出什么
定性结论？
\end{origexamplebox}

\begin{xesolution}
两个导体的净电荷均为正，因此它们的电势均高于无穷远的零电势。取其中电势较高的导体。它是
全空间的最高电势边界，所以其表面电荷密度处处满足 $\sigma\ge0$，不可能出现负电荷区。电势
较低的导体面对高电势导体的一侧，则可能因强烈感应而出现负电荷区，同时在其他区域保留更多
正电荷，使净电荷仍为正。

净电荷较大不等于电势一定较高。孤立导体的电势还取决于尺寸、形状、距离和其他导体的位置，
所以仅由 $Q>q$ 不能断定是 $A$ 还是 $B$ 出现负电荷；能确定的是
\[
  \boxed{\text{两导体不可能同时出现负电荷区}}.
\]

当 $Q=q$ 时，若两导体并不全同且放置不对称，它们的电势仍可不同，上述结论不变。只有在形状、
大小和环境也完全对称时，两者电势相等；此时两块导体共同构成最高电势边界，表面均不会出现
负电荷。

对 $100$ 个净电荷均为正的导体，选电势最高者，它的表面一定处处非负。因此
\[
  \boxed{\text{无论几何怎样，至少有一个导体不含负电荷区。}}
\]
其他较低电势导体是否出现负电荷区仍需由具体几何决定。
\end{xesolution}

\Needspace{0.3\textheight}
\begin{origexamplebox}{例题 31.3【凭空相吸】}
两个金属球本来都不带电。设一次小扰动使某个球内的电荷分布发生变化，进而两个球都因静电感应
而局部带电并相互吸引。这种状态能否自发形成并稳定保持？
\end{origexamplebox}

\begin{xesolution}
不能。两个球都不带净电荷、无外电荷，并取无穷远电势为零。静电平衡时两球各自是等势体，但
电势值尚待确定。若存在宏观电荷分离，就会在球内或球间产生电场；自由电荷立即受力移动，并在
导体电阻、辐射等机制下耗散能量。

最终平衡态可直接由能量判断。静电场能量
\[
  U_E=\frac{\varepsilon_0}{2}\int E^2\dd V\ge0.
\]
在两球均中性且无外场的约束下，$E=0$、两球电势都为零的状态满足全部边界条件，并取得绝对最小
能量 $U_E=0$；由唯一性定理，它也是唯一的静电平衡解。任何宏观电荷分离都对应 $U_E>0$，不能
靠无源小扰动长期维持。

微观热涨落当然会造成瞬时、随机的电荷密度波动，但其平均值为零，且不能按题设方式形成持续的
宏观静电吸引。
\end{xesolution}

\Needspace{0.38\textheight}
\begin{origexamplebox}{例题 31.4【同性相吸】}
在半径为 $R$、净电荷为 $+q$ 的孤立金属球外，沿球心方向距球心 $a>R$ 处放置点电荷 $+q$。
实验中可能出现同性相吸。定量解释这一现象。
\end{origexamplebox}

\begin{xesolution}
先把导体球看成接地球。它对外电荷 $+q$ 的极化作用由
\[
  q_1=-\frac{qR}{a},\qquad r_1=\frac{R^2}{a}
\]
描述。接地球的总感应电荷为 $q_1$，但本题金属球的净电荷必须保持 $+q$，所以还要在球心加入
中心像电荷
\[
  q_0=q-q_1=q\left(1+\frac Ra\right).
\]
中心像电荷提供整体的同号排斥，偏心负像电荷提供近距离极化吸引。外点电荷所受沿球心连线的力为
\begin{align*}
  F&=kq\left[\frac{q_0}{a^2}
  +\frac{q_1}{(a-R^2/a)^2}\right]\\
  &=\boxed{kq^2\left[
  \frac{1+R/a}{a^2}-\frac{Ra}{(a^2-R^2)^2}\right]}.
\end{align*}
规定远离球心为正，$F<0$ 就表示吸引。

令 $x=a/R>1$。临界条件 $F=0$ 化为
\[
  (x-1)^2(x+1)^3=x^4.
\]
其大于 $1$ 的最近实根恰为黄金比
\[
  x=\frac{1+\sqrt5}{2}.
\]
因此
\begin{equation}
  \boxed{R<a<\frac{1+\sqrt5}{2}R\quad\Longrightarrow\quad F<0}.
\end{equation}
点电荷足够靠近球面时，负感应电荷集中在最近处，它产生的吸引随间隙减小而迅速增大，超过金属
球其余正电荷的排斥，于是出现“净电荷同号、合力相吸”。

\isolatedSphereDiagram
\end{xesolution}

\Needspace{0.42\textheight}
\begin{origexamplebox}{例题 31.5【球内屏蔽】}
半径为 $R$ 的接地金属球壳内放置点电荷 $-q$，它距球心为 $x<R$。求：
\begin{enumerate}[label=(\arabic*),leftmargin=2.2em]
  \item 空间中的电场分布；
  \item 感应电荷的面密度分布；
  \item 若球壳不接地而保持净电荷为零，结果怎样变化；
  \item 为什么外部电荷不能在空腔内产生电场，而内部电荷却可能在球壳外产生电场？
\end{enumerate}
\end{origexamplebox}

\begin{xesolution}
取 $z$ 轴沿球心指向真实电荷，真实电荷位于 $z=x$。在球外同一轴线上放置像电荷
\[
  q'=+\frac{qR}{x},\qquad z'=\frac{R^2}{x}.
\]
它位于所求的球内区域之外，不会在球内引入新的奇点。

\textbf{(1) 接地时的电势和电场。}对 $r<R$，
\begin{equation}
  \varphi(\vect r)=k\left[
  -\frac q{|\vect r-x\uvec z|}
  +\frac{qR/x}{|\vect r-(R^2/x)\uvec z|}\right],
  \qquad \vect E=-\nabla\varphi.
\end{equation}
在 $r=R$ 上两项相消。球壳金属内部电场为零；球壳外边界电势为零且无外部电荷，所以球外
$\varphi=0$、$\vect E=0$。

\textbf{(2) 内表面感应电荷。}设 $\theta$ 为球面半径与 $+z$ 轴的夹角。由球内电势沿法线求导，
得到内表面面电荷密度
\begin{equation}
  \boxed{\sigma_{\rm in}(\theta)=
  \frac{q(R^2-x^2)}{4\pi R
  (R^2+x^2-2Rx\cos\theta)^{3/2}}}.
\end{equation}
它处处为正，并在最靠近 $-q$ 的位置 $\theta=0$ 最大。积分结果为
$Q_{\rm in}=+q$。接地时外表面没有电荷。

\textbf{(3) 球壳孤立且净电荷为零。}内表面仍必须感应 $+q$，因为在金属内部取高斯面，包围的
真实电荷与内表面电荷之和必须为零。为了使球壳自身净电荷为零，外表面还要带 $-q$。球形外表面
为等势面，外部又没有其他电荷，所以这部分电荷均匀分布：
\[
  \boxed{\sigma_{\rm out}=-\frac{q}{4\pi R^2}},
  \qquad
  \boxed{\varphi_{\rm out}(r)=-\frac{kq}{r}}.
\]
球内的电场与接地情形相同，电势整体再加常量 $-kq/R$；加常量不改变电场。

\textbf{(4)} 外电荷作用于闭合、无电荷空腔时，空腔边界是常电势，腔内无奇点，唯一解为常量，
所以腔内场为零。内部电荷则在腔内引入了 Poisson 方程的奇点，内表面必须出现等量异号感应
电荷；若导体孤立且净电荷固定，电荷守恒又迫使外表面出现补偿电荷，从而产生外部电场。若球壳
接地，补偿电荷可流入大地，外部场才重新为零。

\innerSphereDiagram
\end{xesolution}

\Needspace{0.34\textheight}
\begin{origexamplebox}{例题 31.6【能量】}
计算前述接地导体球外点电荷系统，以及例题 31.5 中接地球壳内点电荷系统的静电势能。
\end{origexamplebox}

\begin{xesolution}
设接地球外的点电荷为 $Q$，距球心 $L>R$。像电荷为 $-QR/L$，它与真实电荷的距离为
$(L^2-R^2)/L$。真实电荷处的感应电势为
\[
  \varphi_{\rm ind}(Q)
  =-\frac{kQR}{L^2-R^2}.
\]
导体表面电荷是随 $Q$ 从零逐渐建立的，故系统能量不是 $Q\varphi_{\rm ind}$，而是
\begin{equation}
  \boxed{U_{\rm out}=\frac12Q\varphi_{\rm ind}(Q)
  =-\frac{kQ^2R}{2(L^2-R^2)}}.
\end{equation}

例题 31.5 中真实电荷为 $-q$，距球心 $x<R$；像电荷为 $+qR/x$，两者距离为
$(R^2-x^2)/x$。因此
\[
  \varphi_{\rm ind}(-q)=\frac{kqR}{R^2-x^2},
\]
从而
\begin{equation}
  \boxed{U_{\rm in}=\frac12(-q)\varphi_{\rm ind}(-q)
  =-\frac{kq^2R}{2(R^2-x^2)}}.
\end{equation}
两式都为负，且点电荷越靠近球面，能量越低。直接对能量关于 $L$ 或 $x$ 求负梯度，可复核像
电荷给出的吸引力。
\end{xesolution}

\Needspace{0.38\textheight}
\begin{origexamplebox}{例题 31.7【电像法】}
两个相互正交的半无限大接地金属板围成空间的四分之一。以交线为 $z$ 轴，两板分别位于
$x=0$、$y=0$，在第一象限的 $(a,b,0)$ 处放置点电荷 $+q$。求 $x>0,y>0$ 区域的电场。
\end{origexamplebox}

\begin{xesolution}
仅对 $x=0$ 镜像会破坏 $y=0$ 上的零电势，仅对 $y=0$ 镜像也一样。必须同时放置三个像电荷：
\[
  -q\text{ 位于 }(-a,b,0),\qquad
  -q\text{ 位于 }(a,-b,0),\qquad
  +q\text{ 位于 }(-a,-b,0).
\]
关于任一接地平面对称的电荷都成等量异号配对，所以两个平面上的电势同时为零。

记
\[
  \vect r_1=(a,b,0),\quad \vect r_2=(-a,b,0),\quad
  \vect r_3=(a,-b,0),\quad \vect r_4=(-a,-b,0).
\]
第一象限中的电势与电场为
\begin{align*}
  \varphi(\vect r)&=kq\left(
  \frac1{|\vect r-\vect r_1|}-\frac1{|\vect r-\vect r_2|}
  -\frac1{|\vect r-\vect r_3|}+\frac1{|\vect r-\vect r_4|}\right),\\
  \boxed{\vect E(\vect r)}&=\boxed{kq\left[
  \frac{\vect r-\vect r_1}{|\vect r-\vect r_1|^3}
  -\frac{\vect r-\vect r_2}{|\vect r-\vect r_2|^3}
  -\frac{\vect r-\vect r_3}{|\vect r-\vect r_3|^3}
  +\frac{\vect r-\vect r_4}{|\vect r-\vect r_4|^3}
  \right]}.
\end{align*}
三个像电荷都位于真实求解区域之外，故该构造满足唯一性定理的全部条件。

\quarterPlaneDiagram
\end{xesolution}

\Needspace{0.42\textheight}
\begin{origexamplebox}{例题 31.8【半球和平面】}
一个无限大接地平面向上凸起半径为 $R$ 的半球，半球球心位于平面内。在球心正上方
$a>R$ 处放置点电荷 $+q$，求导体外的电场。若半球改为向下凹入平面，为什么不能用同样的
有限个像电荷处理？
\end{origexamplebox}

\begin{xesolution}
取平面为 $z=0$，半球球心为原点。真实电荷 $+q$ 位于 $z=a$。为同时满足平面与半球面上的
零电势，放置三枚像电荷：
\begin{center}
\begin{tabular}{c c}
\toprule
位置 & 电荷量\\
\midrule
$z=-a$ & $-q$\\
$z=R^2/a$ & $-qR/a$\\
$z=-R^2/a$ & $+qR/a$\\
\bottomrule
\end{tabular}
\end{center}
第一枚是 $+q$ 关于接地平面的镜像；第二枚是 $+q$ 关于接地球面的 Kelvin 像；第三枚同时是
前两枚分别关于球面、平面的配对像。

用柱坐标 $(\rho,z)$，在 $z>0$ 且 $\rho^2+z^2>R^2$ 的真实区域内，
\begin{align*}
  \varphi(\rho,z)=kq\bigg[&
  \frac1{\sqrt{\rho^2+(z-a)^2}}
  -\frac1{\sqrt{\rho^2+(z+a)^2}}\\
  &-\frac{R/a}{\sqrt{\rho^2+(z-R^2/a)^2}}
  +\frac{R/a}{\sqrt{\rho^2+(z+R^2/a)^2}}
  \bigg],
\end{align*}
\begin{equation}
  \boxed{\vect E=-\nabla\varphi}.
\end{equation}
在平面上，关于 $z=0$ 的两对电荷分别相消；在半球面上，每个外电荷与它的 Kelvin 像分别
相消，所以两类导体边界都满足 $\varphi=0$。

若半球向下凹入，原来的 $z=R^2/a$ 像点会落入真实可达的凹腔区域，成为不允许出现的虚假奇点。
继续作平面镜像和球面反演会不断产生新的像，有限四电荷构造不再闭合。因此不能照搬上面的有限
像电荷解；通常需采用无穷像电荷级数、分离变量或数值边界法。

\hemisphereImageDiagram
\end{xesolution}

\Needspace{0.4\textheight}
\begin{origexamplebox}{例题 31.9【能量相关】}
无穷大接地金属板的一侧，距板 $a$ 处放置点电荷 $+q$。
\begin{enumerate}[label=(\arabic*),leftmargin=2.2em]
  \item 系统静电势能为多大；
  \item 缓慢将点电荷拉到距板 $2a$ 处，外力做功多少；
  \item 瞬间将点电荷拉到距板 $2a$ 处，外力做功多少；
  \item 后两问的结果为什么不同，差值对应的能量去了哪里？
\end{enumerate}
\end{origexamplebox}

\begin{xesolution}
在金属板另一侧对称位置放置像电荷 $-q$，真实电荷与像电荷的距离为 $2a$。

\textbf{(1)} 真实电荷处的感应电势为 $-kq/(2a)$。系统能量是其一半：
\begin{equation}
  \boxed{U(a)=\frac12q\left(-\frac{kq}{2a}\right)
  =-\frac{kq^2}{4a}}.
\end{equation}

\textbf{(2) 缓慢移动。}每一时刻金属板上的感应电荷都来得及调整到新的静电平衡，所以外力功
等于平衡态能量之差：
\begin{equation}
  \boxed{W_{\rm slow}=U(2a)-U(a)
  =-\frac{kq^2}{8a}+\frac{kq^2}{4a}
  =\frac{kq^2}{8a}}.
\end{equation}

\textbf{(3) 瞬间移动。}在极短移动时间内，板上的感应电荷分布可近似看作仍保持初态分布。初态
感应电荷在板上方产生的电势，等效于固定在 $z=-a$ 的像电荷 $-q$。真实电荷从 $z=a$ 移到
$z=2a$ 时，外力功为
\begin{align*}
  W_{\rm fast}
  &=q\left[\varphi_{\rm old}(2a)-\varphi_{\rm old}(a)\right]\\
  &=q\left(-\frac{kq}{3a}+\frac{kq}{2a}\right)
  =\boxed{\frac{kq^2}{6a}}.
\end{align*}
这里不能再乘 $1/2$：移动阶段被拉动的是完整的真实电荷，而旧感应电荷暂时固定不变。

\textbf{(4)} 瞬移结束后，初态感应电荷分布并不是距离 $2a$ 时的平衡分布。导体内出现短暂电流，
表面电荷重新分布，并伴随焦耳热和少量电磁辐射。两种外力功之差
\begin{equation}
  \boxed{W_{\rm fast}-W_{\rm slow}=\frac{kq^2}{24a}}
\end{equation}
正是随后驰豫过程中耗散掉的能量。慢速过程始终近似静电平衡，没有这部分额外耗散。

\groundPlaneEnergyDiagram
\end{xesolution}

\Needspace{0.4\textheight}
\begin{origexamplebox}{例题 31.10【匀强电场中的导体球】}
半径为 $R$ 的中性导体球置于匀强电场 $\vect E_0$ 中，求空间中的电场分布。
\end{origexamplebox}

\begin{xesolution}
取 $z$ 轴沿 $\vect E_0$，以无穷远处的匀强场电势
$\varphi_0=-E_0r\cos\theta$ 为背景。轴对称、球外无电荷的 Laplace 方程通解中，满足无穷远
条件且只含一阶角向项的形式为
\[
  \varphi_{\rm out}=\left(-E_0r+\frac{A}{r^2}\right)\cos\theta.
\]
导体球面必须是等势面。因球中性且问题关于赤道反对称，可把球面电势常量取为零。代入 $r=R$
得 $A=E_0R^3$，故
\begin{equation}
  \boxed{\varphi(r,\theta)=
  \begin{cases}
  0,&r\le R,\\[3pt]
  -E_0\left(r-\dfrac{R^3}{r^2}\right)\cos\theta,&r>R.
  \end{cases}}
\end{equation}
球外电场的球坐标分量为
\begin{align*}
  \boxed{E_r}&=\boxed{E_0\left(1+2\frac{R^3}{r^3}\right)\cos\theta},\\
  \boxed{E_\theta}&=\boxed{-E_0\left(1-\frac{R^3}{r^3}\right)\sin\theta},
  \qquad E_\phi=0.
\end{align*}
在球面上 $E_\theta(R)=0$，符合导体表面无切向场；法向场
$E_r(R)=3E_0\cos\theta$，所以
\begin{equation}
  \boxed{\sigma(\theta)=3\varepsilon_0E_0\cos\theta}.
\end{equation}
右半球带正电，左半球带负电，总感应电荷为零。附加项
$E_0R^3\cos\theta/r^2$ 等效于球心处的偶极子，其偶极矩为
\begin{equation}
  \boxed{\vect p=4\pi\varepsilon_0R^3\vect E_0}.
\end{equation}

场线在球面处只能沿法向进入或离开，不能穿过导体内部。图 31.11 中靠近两极的场线更密，正好
对应 $|\sigma|$ 在 $\theta=0,\pi$ 处最大；赤道上 $\cos\theta=0$，表面感应电荷密度为零，场线
也与球面相切前后平滑转向。图形中的正、负号表示感应电荷重新分布，并不表示导体球获得了净电荷。

在离球很远的区域 $r\gg R$，附加电势按 $r^{-2}$ 衰减，外场趋近原来的匀强场；在球面附近，
偶极修正与背景场同量级，才显著改变场线。沿电场轴，球面外法向场达到 $3E_0$；沿赤道方向，
法向场为零。用这些轴线极限检查解析式，既能确认正负号，也能把公式与图中的疏密变化对应起来。

若外加场方向反转，感应电荷正负区域同时交换，而球内仍保持零场；若球具有额外净电荷，则还要
在线性叠加中加入球对称的 $Q/(4\pi\varepsilon_0r)$ 项。前者改变偶极项方向，后者增加单极项，
两种效应在图像和远场衰减阶数上可以清楚地区分。

\uniformFieldSphereDiagram

还可以直接检验球的中性条件。把右半球的感应电荷积分：
\[
 Q_+=\int_0^{2\pi}\!\int_0^{\pi/2}
 3\varepsilon_0E_0\cos\theta\,R^2\sin\theta\,\dd\theta\dd\phi
 =3\pi\varepsilon_0E_0R^2,
\]
左半球给出 $Q_-=-Q_+$，故总电荷严格为零。尽管净电荷为零，正、负电荷中心分离仍产生非零
偶极矩；这正是“中性”与“没有电学响应”之间的区别。

球的极化率由 $\vect p=\alpha\vect E_0$ 定义，因此孤立导体球有
\[
 \boxed{\alpha=4\pi\varepsilon_0R^3}.
\]
在均匀场中，诱导偶极子的合力为零；若外场缓慢随空间变化，靠近强场一侧的表面电荷受力更大，
中性导体仍会被吸向场强较大的区域。其低阶受力可写成
$\vec F\simeq\tfrac12\alpha\nabla(E^2)$，其中因子 $1/2$ 来自感应偶极矩随外场建立的过程。

这也解释了为何中性小纸屑会被带电体吸引：纸屑先被外场极化，近端异号电荷受到的吸引强于远端
同号电荷受到的排斥。若把外场近似为严格均匀，只有力矩或内部应力而没有净平动力；必须保留场的
空间不均匀性，才能得到净吸引。图 31.11 给出的是均匀场基准解，非均匀场可在小球尺度远小于
外场变化尺度时，把局部场代入上述偶极近似。
\end{xesolution}
